\documentclass[11pt,a4paper]{article}
\usepackage[utf8]{inputenc}
\usepackage[japanese]{babel}
\usepackage{amsmath,amssymb}
\usepackage{bm}
\usepackage{geometry}

\geometry{margin=25mm}
\title{宇宙物理学 レポート課題}
\author{大阪公立大学理学研究科物理学専攻 BHB26052 河村太星}
\date{\today}

\begin{document}

\maketitle

\section{銀河(Galaxy)の物理的性質と分類}
銀河は、恒星、星間物質(ガスとダスト)、そして目に見えないダークマター(暗黒物質)が重力によって束縛された大規模な天体システムです。
その質量は$10^7 \, M_\odot$(矮小銀河)から$10^{12} \, M_\odot$(巨大楕円銀河)にまで及びます($M_\odot$は太陽質量)。

\subsection{A. ハッブル分類(Hubble Sequence)と形態学}
銀河はそのビジュアル(形態)から、主に以下の4つに分類されます。
\begin{enumerate}
    \item \textbf{楕円銀河（E0〜E7）:} 球形または楕円体状の銀河。星間ガスをほとんど含まず、星形成活動はほぼ終息しています。
    老いた赤い星(種族II)が主成分です。
    \item \textbf{渦巻銀河（Sa〜Sc） / 棒渦巻銀河（SBa〜SBc）:} 明確な「ディスク(銀河円盤)」と「バルジ(中心膨らみ)」を持ちます。
    円盤部には豊富な星間分子雲が存在し、現在も活発な星形成が行われています。
    \item \textbf{レンズ状銀河（S0）:} 円盤はあるものの、渦巻き腕が見られない、楕円銀河と渦巻銀河の中間的な性質を持つ銀河です。
    \item \textbf{不規則銀河（Irr）:} 明確な対称構造や中心核を持たない銀河(例：大マゼラン雲)。ガスが豊富で若い星が多く存在します。
\end{enumerate}

\subsection{B. 暗黒物質(ダークマター)の存在}
銀河の質量分布を調べた際、可視光で見える星やガスの質量だけでは、銀河の回転速度(特に外縁部)を説明できません。外縁部でも回転速度が
落ちない「銀河の回転曲線問題」から、銀河の総質量の約80〜90\%は、電磁波で観測できない\textbf{ダークマター}のハロー(球状の巨大な包み)で
構成されていることが確実視されています。

\section{恒星(Star)の構造と物理機構}
恒星は、中心部での核融合エネルギーによる外向きの熱圧力と、自らの質量による内向きの重力が完全に釣り合った\textbf{「流体静力学平衡
（Hydrostatic Equilibrium）」}の動的状態にあるガス球です。

\subsection{A. 基本的な連続方程式}
恒星の内部構造を記述する際、以下の基礎方程式(流体静力学平衡と質量保存)が根幹となります。
\begin{equation}
    \frac{dP(r)}{dr} = -\frac{G M(r) \rho(r)}{r^2}
\end{equation}

\begin{equation}
    \frac{dM(r)}{dr} = 4\pi r^2 \rho(r)
\end{equation}
ここで、$r$は中心からの距離、$P(r)$は圧力、$\rho(r)$は密度、$M(r)$は半径$r$内部に含まれる質量、 $G$は万有引力定数です。

\subsection{B. 中心エネルギー源(核融合反応)}
主系列星の段階では、中心温度が$\sim 10^7 \, \text{K}$を超えることで水素核融合($4\text{H} \rightarrow \text{}^4\text{He} + \Delta E$)が
駆動します。
\begin{itemize}
    \item \textbf{p-pチェイン（陽子-陽子連鎖反応）:} 太陽質量程度かそれ以下の比較的低温な星で優勢。
    \item \textbf{CNOサイクル:} 太陽より重い星(中心温度がさらに高い星)で優勢。炭素・窒素・酸素を触媒として水素がヘリウムに変わるため、
    エネルギー生成率の温度依存性が非常に高い($\epsilon \propto T^{16}$以上)。
\end{itemize}

\section{恒星の進化経路と質量の依存性}
恒星の一生は、ヘルツシュプルング・ラッセル図(HR図)上での移動として表現されます。進化の分岐点は、電子の縮退圧がいつ、どの段階で
重力崩壊を食い止めるか(あるいは超えるか)という\textbf{質量}の大きさに依存します。

\subsection{A. 低・中質量星($\lesssim 8 \, M_\odot$)の進化}
\begin{enumerate}
    \item \textbf{主系列段階:} 中心で水素が燃焼。全寿命の約90\%をここで過ごします。
    \item \textbf{赤色巨星枝（RGB）:} 中心核の水素が枯渇してヘリウムの核が収換。その周囲の殻(シェル)で水素燃焼が始まり、外層が数倍〜数百倍に膨張。
    \item \textbf{水平枝（HB）/漸近巨星枝（AGB）:} 中心温度が約1億Kに達すると、ヘリウムが炭素や酸素に変わる核融合(トリプルアルファ反応)が始まります。
    \item \textbf{最期:} 外層ガスを宇宙空間に穏やかに放出して\textbf{惑星状星雲}となり、中心に残された炭素・酸素の核は、電子縮退圧で支えられた
    超高密度の\textbf{白色矮星（White Dwarf）}として冷えていきます。
\end{enumerate}

\subsection{B. 大質量星($\gtrsim 8 \, M_\odot$)の進化}
\begin{enumerate}
    \item 中心温度が極めて高くなるため、ヘリウム燃焼以降も、炭素(C)$\rightarrow$ネオン(Ne)$\rightarrow$酸素(O)$\rightarrow$ケイ素(Si)と、
    より重い元素の核融合が連鎖的に進みます。内部は玉ねぎの皮のような層状構造になります。
    \item 中心に最も安定な元素である\textbf{鉄（Fe）の核}が形成されると、それ以上の核融合反応はエネルギー吸収反応となるため、エネルギー生成がストップします。
    \item \textbf{重力崩壊と超新星爆発:} 鉄の核の質量が、電子縮退圧の限界値である\textbf{「チャンドラセカール限界（約$1.4 \, M_\odot$）」}を超えると、
    核が一気に重力崩壊を起こします。電子と陽子が合体して中性子に変わる際の凄まじい衝撃波とニュートリノの放出により、外層が吹き飛ぶ\textbf{「II型超新星爆発」}が
    発生します。
    \item 残骸の質量により、中心には中性子の縮退圧で支えられた\textbf{中性子星（Neutron Star）}、あるいは重力が完全に打ち勝った\textbf{ブラックホール
    （Black Hole）}が形成されます。
\end{enumerate}

\section{恒星の表面温度を決定する観測的手法}
天文学において「星の温度」を直接触って測ることは不可能なため、地球に届く電磁波の情報から物理量を逆算します。主に用いられる指標は\textbf{「実効温度
（Effective Temperature: $T_{\text{eff}}$）」}です。これは、恒星と同じ半径・同じ全光度を持つ黒体の温度として定義されます。

\subsection{A. ウィーンの変位則(黒体放射近似)}
恒星の連続光スペクトルを黒体放射のプランクトン関数$B_\nu(T)$で近似します。放射強度が最大となるピーク波長$\lambda_{\text{max}}$と温度$T$の間には、
以下の反比例の関係が成り立ちます。
\begin{equation}
    \lambda_{\text{max}} = \frac{b}{T} \quad (b \approx 2.8978 \times 10^{-3} \, \text{m} \cdot \text{K})
\end{equation}
このため、高解像度の分光でなくても、どの波長域が最も明るいか(青いか、赤いか)を捉えることで大まかな温度が分かります。

\subsection{B. 光度測定と色指数(Color Index)}
より実用的な方法として、特定の波長透過特性を持つフィルターで星の明るさ(等級)を測定します。標準的なジョンソン・クズンズ(UBVRI)システムでは、
青色(Bバンド：$\sim 440 \, \text{nm}$)と可視・緑色(Vバンド：$\sim 550 \, \text{nm}$)の等級の差である\textbf{$B-V$（色指数）}が多用されます。
\begin{itemize}
    \item 等級は「数値が小さいほど明るい」ため、高温の星(青い光を多く出す)は$B$等級が小さくなり、$B-V$はマイナス、あるいは小さな値になります。
    \item この色指数と実効温度は1対1に対応するため、多数の星の表面温度を効率よくカタログ化する際に強力な手法です。
\end{itemize}

\subsection{C. 分光分類とサハの電離公式(Saha Equation)}
最も厳密な表面温度の決定は、スペクトル中の\textbf{「吸収線（元素の指紋）」}の解析に基づきます。星の表面温度が変わると、大気を構成する原子の電子のエネルギー状態
(励起状態)や、電子が剥ぎ取られる割合(電離度)が劇的に変化します。

これを物理的に記述するのが\textbf{ボルツマン公式}と\textbf{サハの電離公式}です。
\begin{equation}
    \frac{N_{j+1}}{N_j} = \frac{2 k_B T}{P_e} \frac{Z_{j+1}}{Z_j} \left(\frac{2\pi m_e k_B T}{h^2}\right)^{3/2} \exp\left(-\frac{\chi_j}{k_B T}\right)
\end{equation}

\begin{itemize}
    \item $N_j$: 電離段階$j$にある原子の数密度
    \item $P_e$: 電子圧
    \item $\chi_j$: 電離エネルギー
    \item $Z$: 分配関数、$m_e$: 電子質量、$h$: プランク定数、$k_B$: ボルツマン定数
\end{itemize}

この数式が示す通り、温度$T$が変化すると特定の電離状態にある原子の数がピークを迎えます。
例えば、水素のバルマー線(主量子数$n=2$からの遷移)は、温度が低すぎると電子が基底状態($n=1$)に留まるため見えず、逆に温度が高すぎると
水素がすべて電離して(HII)電子を失うため見えなくなります。結果として、バルマー線は\textbf{約10,000 K（A型星）で最も強く現れる}という
性質を持ちます。
天文学では、この吸収線のパターンから恒星を\textbf{O, B, A, F, G, K, M}(高温$\rightarrow$低温)のスペクトル型に分類(ハーバード分類)し、
大気モデルをフィッティングさせることで、数パーセント以下の精度で表面温度$T_{\text{eff}}$を算出しています。

\end{document}